\section{Methodology}

Our three-component system addresses distinct failure modes in current dataset deprecation practice: automated health metrics for candidate detection (failure mode 1), context-aware successor graphs for task-specific recommendations (failure mode 2), and load-time instrumentation for adoption tracking (failure mode 3). Each component targets a separate infrastructure gap; their synergistic integration enables formal deprecation mechanisms validated through five sub-hypothesis experiments.

\subsection{System Architecture}

The system integrates three layers:

\begin{enumerate}
\item \textbf{Health Metrics Layer:} Computes deprecation signals from repository metadata (usage velocity, successor emergence, issue accumulation) without manual curator intervention

\item \textbf{Context-Aware Graph Layer:} Models task-conditional successor relationships (e.g., ImageNet $\rightarrow$ ImageNet-v2 for robustness vs ImageNet-21k for pretraining) inferred from dataset card citations and usage patterns

\item \textbf{Instrumentation Layer:} Wraps dataset loaders with SHA256-based content hashing and async telemetry queues to track adoption events with minimal overhead
\end{enumerate}

Design rationale: Each component solves one failure mode independently but creates synergy through integration---health metrics surface candidates, context graphs personalize recommendations, instrumentation measures outcomes.

\subsection{Component 1: Automated Health Metrics}

Health metrics adapt software package deprecation patterns (NPM, PyPI) to ML's usage-driven lifecycle. We compute weighted scores combining three signals:

\begin{itemize}
\item \textbf{Usage Velocity:} Download rate decline (velocity < 0.3 flags stagnation)
\item \textbf{Successor Emergence:} Alternative dataset availability (emergence > 3 indicates replacements exist)
\item \textbf{Issue Ratio:} Maintenance burden (high issue-to-download ratio signals abandonment)
\end{itemize}

\textbf{Weighted Score:} $\text{score} = 2.0 \times \text{velocity} + 1.5 \times \text{emergence} + 1.0 \times \text{issue\_ratio}$

Velocity weighted highest because declining usage is the strongest deprecation predictor; emergence captures successor availability; issue ratio indicates maintainer burden. Threshold $\geq$2.5 flags deprecation candidates.

\textbf{Design Decision:} Weighted scoring over binary rules provides tunable precision-recall tradeoff. Alternatives considered: uniform weights (lower precision), API-based metrics (blocked by integration complexity), maintainer-only annotations (doesn't scale to 60k+ datasets).

\subsection{Component 2: Context-Aware Successor Graphs}

Successor graphs model task-conditional replacement paths rather than linear version succession. Edges represent deprecation relationships labeled by task context (e.g., \texttt{ImageNet --[robustness\_eval]--> ImageNet-v2}, \texttt{ImageNet --[pretraining]--> ImageNet-21k}).

\textbf{Context Inference:} Pattern-based matching on Python import history and dataset card citations:
\begin{itemize}
\item Exact match: \texttt{from datasets import load\_dataset("squad")} $\rightarrow$ task="question\_answering"
\item Substring: \texttt{"robustness"} in notebook $\rightarrow$ task="robustness\_eval"
\item Semantic similarity: Dataset card keywords match known task taxonomies
\end{itemize}

\textbf{Graph Construction:} NetworkX DiGraph with datasets as nodes, successor edges inferred from dataset card citations (\texttt{"recommended\_successor: dataset\_name"} fields). Users can override inferred context via explicit task parameter.

\textbf{Design Decision:} Pattern-based inference enables validation on synthetic data; defer embedding-based inference (LLMs) to production testing. Alternatives: manual annotation (doesn't scale), usage co-occurrence (requires longitudinal data unavailable at PoC scale).

\subsection{Component 3: Load-Time Instrumentation}

Decorator-based wrapper for \texttt{datasets.load\_dataset()} adds SHA256 content hashing and async telemetry logging:

\begin{verbatim}
@instrumented_loader
def load_dataset(name, *args, **kwargs):
    # 1. Compute SHA256 hash (lazy caching)
    content_hash = compute_hash(name)
    # 2. Check deprecation status via health metrics
    if is_deprecated(name):
        successor = get_context_aware_successor(name, infer_context())
        warn(f"{name} deprecated. Recommended: {successor}")
    # 3. Log event to async queue (non-blocking)
    telemetry_queue.put({
        "dataset": name,
        "hash": content_hash,
        "successor_shown": successor,
        "timestamp": now()
    })
    # 4. Proceed with normal load
    return original_load_dataset(name, *args, **kwargs)
\end{verbatim}

\textbf{Overhead Mitigation:} SHA256 hashing amortized via lazy evaluation (compute once, cache result); async queue writes prevent blocking I/O.

\textbf{Design Decision:} Decorator pattern enables drop-in replacement without API changes. Alternatives: monkey-patching (fragile), separate CLI tool (adoption barrier), server-side tracking (privacy concerns mitigated via client-side hashing with no user PII).

\subsection{Sub-Hypothesis Decomposition}

We validate mechanisms through five sub-hypotheses with gate protocols:

\begin{itemize}
\item \textbf{h-e1 (MUST\_WORK):} Instrumentation overhead <10\%, capture rate $\geq$95\%
\item \textbf{h-m1 (MUST\_WORK):} Health metrics precision $\geq$60\%, recall $\geq$80\%
\item \textbf{h-m2 (MUST\_WORK):} Context inference accuracy $\geq$70\%, override rate <50\%
\item \textbf{h-m3 (SHOULD\_WORK):} Dependency graph impact completeness $\geq$95\%
\item \textbf{h-m4 (SHOULD\_WORK):} Full-stack telemetry capture $\geq$95\%, overhead <10\%
\end{itemize}

Gate categories (MUST\_WORK vs SHOULD\_WORK) prioritize critical components (detection, recommendation, measurement) over supporting infrastructure (dependency analysis).

\subsection{Implementation Scope}

Proof-of-concept validation uses:
\begin{itemize}
\item \textbf{Mock data (h-m2, h-m3, h-m4):} Isolates mechanism validation from external dependencies
\item \textbf{Simulated HuggingFace metadata (h-m1):} 1000-dataset corpus with synthetic health signals
\item \textbf{PoC scale (h-e1):} 100 dataset loads to validate overhead/capture feasibility
\end{itemize}

Real HuggingFace API integration and production-scale testing are future work. Mock data enables controlled mechanism validation; generalization requires deployment testing.
